\section{Introduction}

\subsection{Statement of the theorem}

\subsubsection{} We fix a smooth proper curve $X$ over an algebraically closed field of characteristic zero, as well as a square root $\omega_X^{1/2}$ of the canonical line bundle.

\subsubsection{} The geometric Langlands conjecture \cite{ss_indcoh} predicts an equivalence of categories
    \[\GLL_G : \DMod(\Bun_G) \simto \ICoh_\Nilp(\LS_\LG).\]
Let us review some of the compatibilities expected from this functor.

\subsubsection{Duality} The category appearing on the spectral side of the Langlands equivalence has a self-duality given by Serre duality
    \[
    \DD : \ICoh_\Nilp(\LS_\LG) \simto \ICoh_\Nilp(\LS_\LG)^\vee
    \qquad
    \FF \mapsto \Gamma (X, \FF \shimes -)
    \]
The work \cite{strange} constructs a self duality
    \[
    \PsId_! : \DMod(\Bun_G)^\vee \simto \DMod(\Bun_G)
    \]
which is expected to be the automorphic counterpart of Serre duality in the sense that the following diagram commutes up to a Cartan involution and cohomological shift.
\begin{equation}
\label{eqn:Serre_PsId_no_shift}
\begin{tikzcd}[
  column sep={9em,between origins},
  row sep={4em,between origins}
]
  \DMod(\Bun_G) \arrow[r, "\GLL_G"] \arrow[dr, phantom, "\tau \cdot {[?]}"] & \ICoh_\Nilp(\LS_\LG) \arrow[d, "\DD"] \\
  \DMod(\Bun_G)^\vee \arrow[u, "\PsId_!"] & \ICoh_\Nilp(\LS_\LG)^\vee \arrow[l, "\GLL_G^\vee"]
\end{tikzcd}
\end{equation}
By definition, $\PsId_!$ is given by the object
    \[
    \Delta_! \cdot k_{\Bun_G} \in
    \DMod(\Bun_G \times \Bun_G)
    \overset{\sim}{\leftarrow} \DMod(\Bun_G) \otimes \DMod(\Bun_G),
    \]
where $\Delta : \Bun_G \to \Bun_G \times \Bun_G$ is the diagonal.

\subsubsection{Whittaker coefficients} Let $\rhov$ be the half-sum of the positive coroots. By using the chosen square root of the canonical bundle, we obtain a point of $\Bun_T$ denoted $\omega^\rhov$ and define
    \[
    \Bun_N = \{\omega^\rhov\} \underset{\Bun_T}{\times} \Bun_B.
    \]
This stack is equipped with canonical maps
    \[
    \begin{tikzcd}[
  column sep={3em,between origins},
  row sep={3em,between origins}
]
      & \Bun_N \arrow[dl, "\psi" above left] \arrow[dr, "\pi"] & \\
      \Ga & & \Bun_G
    \end{tikzcd}
    \]
Now let $\AS$ denote the exponential D-module on $\Ga$, cohomologically shifted to be a character sheaf with respect to $!$-convolution. We write
    \[\Vect \simto \Whit(\Bun_N) \overset{\oblv}{\longrightarrow} \DMod(\Bun_N)\]
for the full subcategory generated by the object $\WW_\vac = \ev^! \cdot \AS$. The inclusion admits a right adjoint
    \[
    \DMod(\Bun_N) \overset{\Av_*}{\longrightarrow} \Whit(\Bun_N) \simeq \Vect,
    \]
and the functor of Whittaker coefficients is
    \[
    \coeff_* = \Av_* \cdot \pi^!.
    \]
It is supposed to match the functor of pushforward along $f : \LS_\LG \to \pt$.
\begin{equation}
\label{eqn:L_G_coeff_Poinc}
\begin{tikzcd}[
  column sep={8em,between origins},
  row sep={4em,between origins}
]
  \DMod(\Bun_G) \arrow[d, "\coeff_*" left] \arrow[r, "\GLL_G"] & \ICoh_\Nilp(\LS_\LG) \arrow[d, "f_*"] \\
  \Whit(\Bun_N) \arrow[r, "\sim"] & \QCoh(\pt)
\end{tikzcd}
\qquad
\begin{tikzcd}[
  column sep={8em,between origins},
  row sep={4em,between origins}
]
    \Whit(\Bun_N) \arrow[r, "\sim"] \arrow[d, "\Poinc_!" left] & \QCoh(\pt) \arrow[d, "f^*"] \\ 
    \DMod(\Bun_G) \arrow[r, "\GLL_G"] & \ICoh_\Nilp(\LS_\LG)
\end{tikzcd}
\end{equation}
Passing to left adjoints gives the diagram on the right, where
    \[\Poinc_! = \pi_! \cdot \oblv.\]

\subsubsection{} It is a general fact that Serre duality exchanges $f_*$ with $f^!$. However, a special feature of $\LS_\LG$ is that there is a canonical isomorphism
    \[\omega_{\LS_\LG} \simeq \OO_{\LS_\LG} \, [2 \delta_G]\]
where
    \[2 \delta_G = (2g - 2) \cdot \dim(\LG) = 2 \dim \Bun_G.\]
Therefore the dual functor to $f_*$ identifies with a shift of $f^*$.
\[
\begin{tikzcd}[
  column sep={10em,between origins},
  row sep={4em,between origins}
]
    \QCoh(\pt) \arrow[r, "\sim"] \arrow[d, "f^*{[2 \delta_G]}" left] & \QCoh(\pt)^\vee \arrow[d, "(f_*)^\vee"] \\ 
    \ICoh_\Nilp(\LS_\LG) \arrow[r, "\DD"] & \ICoh_\Nilp(\LS_\LG)^\vee
\end{tikzcd}
\]
Let us translate this statement to the automophic side. We will use the Verdier duality pairing
    \[
    \Whit_{-\psi}(\Bun_N) \otimes \Whit_\psi(\Bun_N) \to \Vect
    \qquad
    \langle \WW_1, \WW_2 \rangle = \Gamma(\Bun_N, \WW_1 \shimes \WW_2).
    \]
The resulting statement is our main theorem:

\begin{theorem}
There is a canonical commutative diagram
\[
\label{thm:vac}
\begin{tikzcd}[
  column sep={10em,between origins},
  row sep={4em,between origins}
]
    \Whit_{-\psi}(\Bun_N)^\vee \arrow[d, "\coeff_*^\vee {[2 \delta_G]}" left] \arrow[r] & \Whit_\psi(\Bun_N) \arrow[d, "\Poinc_!"] \\
    \DMod(\Bun_G)^\vee \arrow[r, "\PsId_!"] & \DMod(\Bun_G)
\end{tikzcd}
\]
\end{theorem}

\subsubsection{} As we have just explained, Theorem \ref{thm:vac} is predicted by the geometric Langlands conjecture. However, we expect that it will play an essential r\^ole in the construction of the Langlands equivalence $\GLL_G$ and its attendant compatibilities. 

\subsubsection{Cohomological shifts} Due the ambiguity in (\ref{eqn:Serre_PsId_no_shift}), the above discussion does not actually predict the cohomological shift appearing in our theorem. (One can guess this shift by considering the compatibility expressed by Proposition \ref{proposition:poinc_identification}, which only involves the automorphic side of the Langlands correspondence.) Let us use the theorem statement to determine the shift that should appear in (\ref{eqn:Serre_PsId_no_shift}). The Verdier duality on $\Whit(\Bun_N)$ is related to Serre duality as follows:
\[
\begin{tikzcd}[
  column sep={9em,between origins},
  row sep={4em,between origins}
]
  \Whit_\psi(\Bun_N) \arrow[r, "\sim"] \arrow[d, "\DD{[-2 \delta_N]}" left] & \QCoh(\pt) \arrow[d, "\DD"] \\
  \Whit_{-\psi}(\Bun_N)^\vee \arrow[r, "\sim"] & \QCoh(\pt)^\vee
\end{tikzcd}
\]
The shift appearing here is actually to be expected: it accounts for the difference between Verdier duality and the more natural local Whittaker duality. In any case we obtain:
\[
\begin{tikzcd}[
  column sep={9em,between origins},
  row sep={4em,between origins}
]
  \DMod(\Bun_G) \arrow[r, "\GLL_G"] \arrow[dr, phantom, "\tau"] & \ICoh_\Nilp(\LS_\LG) \arrow[d, "\DD {[2\delta_G]}"] \\
  \DMod(\Bun_G)^\vee \arrow[u, "\PsId_! {[2 \delta_G + 2 \delta_N]}"] & \ICoh_\Nilp(\LS_\LG)^\vee \arrow[l, "\GLL_G^\vee"]
\end{tikzcd}
\]

\subsection{Upgraded version}

\subsubsection{} Our main theorem has a $\Sph_{G,\Ran}$-linear upgrade involving the entire factorization category $\Whit(G)_\Ran$. Let us recall some basic facts about Whittaker categories and Hecke actions before giving the statement.

\subsubsection{Global Whittaker model} Let $x$ be a finite set of points of $X$. We will use the global model for the Whittaker category
    \[
    \Whit(G)_x \simeq \Whit\left(\Bunc_{N,x}^\pol\right) \mon \DMod\left(\Bunc_{N,x}^\pol \right).
    \]
It is naturally a module over the spherical Hecke category $\Sph_{G,x}$, equipped with a pair of $\Sph_{G,x}$-linear adjoint functors
    \[
    \Poinc_! : \Whit(G)_x \to \DMod(\Bun_G)
    \qquad
    \coeff_* : \DMod(\Bun_G) \to \Whit(G)_x.
    \]
Here, as before, $\Poinc_!$ is simply $!$-pushforward of the underlying sheaf to $\Bun_G$. The Verdier pairing
    \[
    \Whit_{-\psi}(G)_x \otimes \Whit_\psi(G)_x \to \Vect
    \]
is a perfect pairing.

\subsubsection{Left and right modules} To formulate the sense in which our dualities are $\Sph_G$-linear, recall that every convolution stack comes with a canonical involution
    \[
    \sigma : X \underset{S}{\times} X \to X \underset{S}{\times} X
    \]
given by swapping the two factors. For example, this involution on $G \simeq \pt \underset{\BB G}{\times} \pt$ is group inversion. We obtain an involutive anti-automorphism
    \[\sigma : \Sph_{G,x} \simto \Sph_{G,x}^\rev.\]
Now let us view $\DMod(\Bun_G)$ as a right $\Sph_{G,x}$-module. Then $\DMod(\Bun_G)^\vee$ is naturally a left $\Sph_{G,x}$-module, but the involution $\sigma$ allows us to identify left and right modules for $\Sph_{G,x}$. Therefore, it makes sense to say that the functor
    \[\PsId_! : \DMod(\Bun_G)^\vee \to \DMod(\Bun_G)\]
is canonically $\Sph_{G,x}$-linear. This is easy to check using the ULA property of Hecke kernels and the fact that $\sigma$ interchanges the geometrically defined left and right $\Sph_{G,x}$ actions on $\DMod(\Bun_G)$. It is tautological that the duality
    \[\Whit_{-\psi}(G)_x^\vee \simeq \Whit_\psi(G)_x\]
is $\Sph_{G,x}$-linear.

\subsubsection{} Allowing the set of points $x$ to vary in families and writing 
    \[\DMod(\Bun_G)_\Ran = \DMod(\Bun_G \times \Ran),\]
we obtain the following variant of the main theorem:

\begin{theorem}
\label{thm:Ran}
There is a canonical $\Sph_{G,\Ran}$-linear commutative diagram
\[
\begin{tikzcd}[
  column sep={10em,between origins},
  row sep={4em,between origins}
]
    \Whit_{-\psi}(G)_\Ran^\vee \arrow[d, "\coeff_*^\vee {[2 \delta_G]}" left] \arrow[r] & \Whit_\psi(G)_\Ran \arrow[d, "\Poinc_!"] \\
    \DMod(\Bun_G)_\Ran^\vee \arrow[r, "\PsId_!"] & \DMod(\Bun_G)_\Ran
\end{tikzcd}
\]
\end{theorem}

\subsubsection{} Let us explain why this theorem is to be expected from the point of view of Langlands duality. In our derivation of the main theorem for the Whittaker vacuum, we actually ignored the Cartan involution without comment. This was acceptable because the Cartan involution sends $\OO_{\LS_\LG}$ to itself. But in what follows we will have to track this involution more closely.

\subsubsection{} Recall that geometric Satake defines a monoidal functor
    \[\Sat_x : \Rep(\LG)_x \to \Sph_{G,x}.\]
The Langlands functor is supposed to intertwine the Hecke action of $\Sph_{G,x}$ with the action of $\Rep(\LG)_x$ on $\ICoh_\Nilp(\LS_\LG)$ induced by the map
    \[\ev_x : \LS_\LG(X) \to \LS_\LG(\mathscr{D}_x) \simeq (\BB \LG)_x\]
and the action of $\QCoh$ on $\ICoh$. Therefore, viewing $\WW_\vac$ as an object of $\Whit(G)_x$ and acting on (\ref{eqn:L_G_coeff_Poinc}), we obtain:
\[
\begin{tikzcd}[
  column sep={8em,between origins},
  row sep={4em,between origins}
]
    \Whit(G)_x \arrow[d, "\Poinc_!"] & \Rep(\LG)_x \arrow[d, "\ev_x^*"] \arrow[l, "\CS" above] \\ 
    \DMod(\Bun_G) \arrow[r, "\GLL_G" below] & \ICoh_\Nilp(\LS_\LG)
\end{tikzcd}
\qquad
\begin{tikzcd}[
  column sep={8em,between origins},
  row sep={4em,between origins}
]
  \DMod(\Bun_G) \arrow[d, "\coeff_*"] \arrow[r, "\GLL_G"] & \ICoh_\Nilp(\LS_\LG) \arrow[d, "\ev_{x,*}"] \\
  \Whit(G)_x & \Rep(\LG)_x \arrow[l, "\CS"]
\end{tikzcd}
\]

\subsubsection{Cartan involution} Let us recall the behavior of the involution $\sigma$ with respect to geometric Satake. To reduce confusion, we will temporarily consider the derived Satake isomorphism
    \[
    \Sph_{G,x} \simeq
    \Sph_{\LG,x}^\spec \simeq \ICoh_\Nilp\left( ( \pt \underset{\Lg}{\times} \pt ) /\LG \right).
    \]
The na\"ive diagram
\[
\begin{tikzcd}[
  column sep={8em,between origins},
  row sep={4em,between origins}
]
    \Sph_{G,x} \arrow[r, "\Sat"] \arrow[d, "\sigma" left] \arrow[dr, phantom, "(\tau)"] & \Sph_{\LG,x}^\spec \arrow[d, "\sigma"] \\
    \Sph_{G,x}^\rev \arrow[r, "\Sat^\rev" below] & \Sph_{\LG,x}^{\spec, \rev}
\end{tikzcd}
\]
of algbera isomorphisms is not commutative: the commutativity is only restored after inserting the Cartan involution. Note that the diagram
\begin{equation}
\label{eqn:Sat_Cartan}
\begin{tikzcd}[
  column sep={7em,between origins},
  row sep={4em,between origins}
]
    \Rep(\LG) \arrow[r] \arrow[d, "\id" left] & \Sph_{\LG,x}^\spec \arrow[d, "\sigma"] \\
    \Rep(\LG)^\rev \arrow[r,] & \Sph_{\LG,x}^{\spec, \rev}
\end{tikzcd}
\end{equation}
is canonically commutative because the horizonal map is given by pushforward along the diagonal
    \[\BB \LG \to \BB\LG \underset{\Lg/\LG}{\times} \BB\LG.\]
Since Serre duality on $\ICoh_\Nilp(\LS_\LG)$ is naturally $\Rep(\LG)$-linear, the Cartan involution in (\ref{eqn:Sat_Cartan}) explains the presence of the Cartan involution in (\ref{eqn:Serre_PsId_no_shift}). Similarly, the Cartan involution also appears in the Whittaker duality:
\[
\begin{tikzcd}[
  column sep={8em,between origins},
  row sep={4em,between origins}
]
  \Whit_\psi(G)_x \arrow[d, "\DD{[-2\delta_N]}" left] \arrow[dr, phantom, "\tau"] & \Rep(\LG)_x \arrow[d] \arrow[l, "\CS" above]\\
  \Whit_{-\psi}(G)_x^\vee \arrow[r, "\CS^\vee"] & \Rep(\LG)_x^\vee
\end{tikzcd}
\]

\subsubsection{} Applying the expected compatibilities to every object appearing Theorem \ref{thm:Ran}, we obtain the diagram
\[
\begin{tikzcd}[
  column sep={10em,between origins},
  row sep={4em,between origins}
]
      \Rep(\LG)_x^\vee \arrow[d, "\ev_{x,*}^\vee" left] & \Rep(\LG)_x \arrow[l] \arrow[d, "\ev_x^*"] \\
      \ICoh_\Nilp(\LS_\LG)^\vee & \ICoh_\Nilp(\LS_\LG) \arrow[l, "\DD{[2\delta_G]}"]
\end{tikzcd}
\]

\subsection{Outline of the proof}

\subsubsection{Notation} Let us write $(-)_\co$ for the Verdier duality $\DD(-)$ and
    \[\PsId_\Whit : \Whit_{-\psi}(\Bun_N)^\vee \to \Whit_\psi(\Bun_N)\]
for its inverse.

\subsubsection{} Our first task will be the construction of a natural transformation
    \[\eta_G : \Poinc_! \cdot \PsId_\Whit \to \PsId_! \cdot \coeff_*^\vee \, [2\delta_G].\]
This problem is actually very concrete: the data of $\eta_G$ is equivalent to a morphism
    \[
    \Poinc_! \cdot \WW_\vac \to \PsId_! \cdot \coeff_*^\vee \cdot \WW_{\vac, \co} \, [2 \delta_G]
    \]
between two objects of $\DMod(\Bun_G)$. We will perform this construction in \S\ref{section:construction}. Once this has been done, we will have to prove that $\eta_G$ is an isomorphism.

\subsubsection{} Our proof that $\eta_G$ is an isomorphism is based on the decomposition of $\DMod(\Bun_G)$ into its cuspidal and Eisenstein parts. Recall that the full subcategory $\DMod(\Bun_G)_\cusp \mon \DMod(\Bun_G)$ of cuspidal objects has the following equivalent characterizations:
\begin{itemize}
    \item It consists of objects $\FF$ which are annihilated by $\CT_{P,*}$ for every proper parabolic $P$.
    \item It consists of objects $\FF$ which are annihilated by $\CT_{P,!}$ for every proper parabolic $P$.
    \item It is the right orthogonal to $\DMod(\Bun_G)_\Eis$.
\end{itemize}
Here the full subcategory $\DMod(\Bun_G)_\Eis \mon \DMod(\Bun_G)$ is characterized by the following equivalent characterizations:
\begin{itemize}
    \item It is the full subcategory generated by the images of $\Eis_{P,!}$ for all proper parabolics $P$.
    \item It is the left orthogonal to $\DMod(\Bun_G)_\cusp$.
\end{itemize}
The category $\DMod(\Bun_G)_\Eis$ contains but is not usually generated by the images of $\Eis_{P,*}$ for all proper parabolics $P$. We will prove Theorem \ref{thm:vac} by establishing the following:

\begin{lemma}
\label{lemma:cofib_eta_Eis}
The cofiber of $\eta_G$ is Eisenstein series.
\end{lemma}

\begin{lemma}
\label{lemma:cofib_eta_cusp}
The cofiber of $\eta_G$ is cuspidal.
\end{lemma}

\subsection{Comparison with the na\"ive pseudo-identity}

\subsubsection{} Recall that there is a na\"ive variant of the pseudo-identity functor, denoted
    \[\PsId_* : \DMod(\Bun_G)^\vee \to \DMod(\Bun_G),\]
corresponding to the object
    \[\Delta_* \cdot \omega_{\Bun_G} \in \DMod(\Bun_G \times \Bun_G).\]
It is related to $\PsId_!$ by a canonical map
    \[
    \PsId_! \, [2 \delta_G + z_G] \to \PsId_*
    \qquad
    z_G = \dim Z_G
    \]
whose cofiber is Eisenstein series.

\subsubsection{} The functor $\Poinc_!$ has a variant
    \[\Poinc_* = \pi_* \cdot \oblv,\]
related to $\Poinc_!$ by a canonical map
\begin{equation}
    \label{eqn:Poinc*!2}
    \Poinc_! \, [z_G] \to \Poinc_*.
\end{equation}
The version of our main theorem with $\PsId_*$ and $\Poinc_*$ is actually trivial:

\begin{proposition}
\label{proposition:trivial_thm}
There is a canonical commutative diagram
\[
\begin{tikzcd}[
  column sep={10em,between origins},
  row sep={4em,between origins}
]
    \Whit_{-\psi}(\Bun_N)^\vee \arrow[d, "\coeff_*^\vee" left] \arrow[r, "\PsId"] & \Whit_\psi(\Bun_N) \arrow[d, "\Poinc_*"] \\
    \DMod(\Bun_G)^\vee \arrow[r, "\PsId_*" below] & \DMod(\Bun_G)
\end{tikzcd}
\]
\end{proposition}

\subsubsection{Proof} Performing base change on the diagram
\[
\begin{tikzcd}[
  column sep={4em,between origins},
  row sep={4em,between origins}
]
    & \Bun_N \arrow[dl] \arrow[dr] & & \\
    \Bun_G \arrow[dr, "\Delta"] & & \Bun_N \times \Bun_G \arrow[dl, "\pi \times \id"] \arrow[dr, "\pi_2"] & \\
    & \Bun_G \times \Bun_G & & \Bun_G
\end{tikzcd}
\]
gives
\[
    \PsId_* \cdot \coeff_*^\vee \cdot \WW_\co
    \simeq \pi_{2,*} \cdot \left\{ \WW \shimes (\pi \times \id)^! \cdot \Delta_* \cdot \omega_{\Bun_G} \right\}
    \simeq \Poinc_* \cdot \WW.
\]

\subsubsection{} Combining $\eta_G$ with the morphism relating $\PsId_!$ to $\PsId_*$, we obtain a map
\begin{equation}
\label{eqn:Poinc*!1}
\begin{split}
  \Poinc_! \, [z_G]
  & \to \PsId_! \cdot \coeff^\vee \cdot \PsId_\Whit^{-1} \, [2 \delta_G + z_G] \\
  & \qquad \qquad \to \PsId_* \cdot \coeff^\vee \cdot \PsId_\Whit^{-1} \simeq \Poinc_*.
\end{split}
\end{equation}
Then Lemma \ref{lemma:cofib_eta_Eis} is equivalent to the assertion that the cofiber of this map is Eisenstein series. We will prove this in \S\ref{section:cofib_Eis} by verifying the following two assertions:

\begin{proposition}
\label{proposition:poinc_identification}
The map (\ref{eqn:Poinc*!1}) canonically identifies with (\ref{eqn:Poinc*!2}).
\end{proposition}

\begin{proposition}
\label{proposition:poinc_principal_series}
The cofiber of (\ref{eqn:Poinc*!2}) is Eisenstein series.
\end{proposition}

\subsubsection{Remark} Combining Theorem \ref{thm:vac} with Proposition \ref{proposition:trivial_thm} gives
    \[
    \Poinc_* \simeq \DL \cdot \Poinc_!
    \]
where
    \[\DL = \PsId_* \cdot \PsId_!^{-1} \, [-2\delta_G]\]
is the Deligne--Lusztig functor. This means that $\Poinc_* \cdot \WW_\vac$ should correspond under Langlands duality to the Steinberg object constructed by D. Beraldo \cite{dario_DL}.

\subsection{Jacquet functors} 

\subsubsection{} In order to show that the cofiber of $\eta_G$ is cuspidal, we need to prove that the morphism
    \[
    \CT_{P, !}^- \cdot \Poinc_! \cdot \WW_\vac 
    \to \CT_{P, !}^- \cdot \PsId_! \cdot \coeff^\vee \cdot \WW_{\vac,\co} \, [2\delta_G]
    \]
is an equivalence for every proper parabolic $P$. We will do this by writing both sides in terms of Poincare series for $M$ and applying induction. This will force us to consider the Ran version of the Whittaker category.

\subsubsection{The LHS} We will construct a Jacquet functor 
  \[\Jac : \Whit(\Bun_N) \to \Whit(M)_\Ran\]
fitting into a commutative square:
\[
\begin{tikzcd}[
  column sep={7.5em,between origins},
  row sep={4em,between origins}
]
    \Whit(G) \arrow[d, "\Poinc_!" left] \arrow[r, "\Jac"] & \Whit(M)_\Ran \arrow[d, "\Poinc_!"] \\
    \DMod(\Bun_G) \arrow[r, "\CT_P^-"] & \DMod(\Bun_M) 
\end{tikzcd}
\]
The right vertical functor appearing in this diagram is the composite
    \[\Whit(M)_\Ran \to \DMod(\Bun_M)_\Ran \to \DMod(\Bun_M),\]
where the first arrow is the family version of Poincare series appearing in Theorem \ref{thm:Ran} for $M$, and the second arrow is integration over the Ran space. The proof that this square is commutative will require geometric input in the form of a ULA property established in \S\ref{section:ULA}.

\subsubsection{The RHS} First recall that
    \[
    \CT^-_{P, !} \cdot \PsId_{G,!} 
    \simeq \PsId_{M,!} \cdot (\Eis_{P,!}^-)^\vee
    \]
because, as objects of $\DMod(\Bun_G \times \Bun_M)$, both identify by base-change with
    \[(p \times q)_! \cdot k_{\Bun_{P^-}}.\]
Then, similarly to the above, we will construct a functor
    \[\Jac_\co : \Whit(M)_\Ran \to \Whit(\Bun_N)\]
equipped with a commutative square:
\[
\begin{tikzcd}[
  column sep={8em,between origins},
  row sep={4em,between origins}
]
  \DMod(\Bun_M) \arrow[r, "\Eis_!^-"] \arrow[d, "\coeff_*" left] & \DMod(\Bun_G) \arrow[d, "\coeff_*"] \\
  \Whit(M)_\Ran \arrow[r, "\Jac_{\co}"] & \Whit(G)
\end{tikzcd}
\]
The functors $\Jac$ and $\Jac_\co$ will be dual to each other in the sense that there a commutative diagram:
\[
\begin{tikzcd}[
  column sep={8em,between origins},
  row sep={4em,between origins}
]
  \Whit(G) \arrow[r, "\Jac"] \arrow[d, "\DD{[2\delta_G]}" left] & \Whit(M)_\Ran \arrow[d, "\DD{[2\delta_M]}"] \\
  \Whit(G)^\vee \arrow[r, "\Jac_{\co}^\vee"] & \Whit(M)_\Ran^\vee 
\end{tikzcd}
\]

\subsubsection{} We organize this discussion by pasting everything into a large diagram.
\begin{equation}
\label{eqn:the_cube}
\begin{tikzcd}[
  column sep={5em,between origins},
  row sep={4em,between origins}
]
  \Whit(M)_\Ran^\vee \arrow[dddd, "\coeff_*^\vee" left] & & & & \Whit(M)_\Ran \arrow[llll, "\DD{[2\delta_M]}" above] \arrow[dddd, "\Poinc_!"] \\
  & \Whit(G)^\vee \arrow[dd, "\coeff_*^\vee" left] \arrow[ul, "\Jac_\co^\vee" above right] & & \Whit(G) \arrow[ll, "\DD{[2\delta_G]}" above] \arrow[dd, "\Poinc_!"] \arrow[ur, "\Jac"] & \\
  & & & & \\
  & \DMod(\Bun_G)^\vee \arrow[dl, "\Eis_{P,!}^{-,\vee}"] \arrow[rr, "\PsId_!"] & & \DMod(\Bun_G) \arrow[dr, "\CT_{P,!}^-" below left]  &  \\
  \DMod(\Bun_M)^\vee \arrow[rrrr, "\PsId_{M,!}" below] & & & & \DMod(\Bun_M)
\end{tikzcd}
\end{equation}
In \S\ref{section:Ran_upgrades}, we shall extend the construction of $\eta_G$ to a $\Sph_{G,\Ran}$-linear morphism
    \[
    \eta_{G,\Ran} : \Poinc_{\Ran, !} \cdot \PsId_{\Ran, !}^\Whit \to \PsId_{G,!} \cdot \coeff_{\Ran, *}^\vee \, [2\delta_G].
    \]
It is compatible in the sense that its evaluation on $\WW_\vac$ agrees with $\eta_G$. Therefore, Theorem \ref{thm:vac} for a given group implies Theorem \ref{thm:Ran} for that same group.

\begin{lemma}
\label{lemma:the_cube_commutes}
Complete (\ref{eqn:the_cube}) into a hexahedron of functors and natural transformations by pasting $\eta_{M,\Ran}$ along the outer edges. Then there exists a 3-morphism rendering the cube commutative.
\end{lemma}

\subsubsection{} We were not able to prove Lemma \ref{lemma:the_cube_commutes} directly. Instead, the proof of the theorem goes through a slightly different statement from which this lemma can be deduced. This is a rather unsatisfactory situation because the resulting 3-morphism is non-canonical.

\subsubsection{} Let us explain how Theorem \ref{thm:vac} would follow by induction on $G$ if we knew Lemma \ref{lemma:the_cube_commutes}. Suppose that $\eta_M$ is an isomorphism for all proper Levi subgroups of $G$. By $\Sph_{M,\Ran}$-linearity, $\eta_{M,\Ran}$ is a an isomorphism for all proper Levis $M$. Then Lemma \ref{lemma:the_cube_commutes} implies that $\CT_{P,!}^- \cdot \eta_G$ is an isomorphism for all proper Levis $M$. This proves Lemma \ref{lemma:cofib_eta_cusp} for $G$. Combining this with Lemma \ref{lemma:cofib_eta_Eis} proves that $\eta_G$ is an isomorphism.

\subsection{Acknowledgements}

\subsubsection{} This problem was suggested by our advisor D. Gaitsgory, and this work would not have existed without his steady guidance and support. We are grateful to S. Raskin for conversations which lead to the argument of \S\ref{section:ULA}. We are also grateful to L. Chen and Y. Fu for many patient explanations and for their encouragement.

\subsubsection{} We thank L. Chen and D. Gaitsgory for their comments on an earlier draft of this work.